\section{Perfect, Whole, and Complete}

Esoterists have historically been prudent about making their teachings known to the general public. Hence, the teachings were disguised in parables, proverbs, myths, symbols, cryptic sayings, and so on. On the one hand, this prevented the teachings from being distorted and misunderstood, if they were presented in a more open way. On the other hand, the teachings were preserved since the stories and symbols would be pass on, even by those who could not fully comprehend them.

In the past, books were rare and expensive, so teachings had to be passed down from mouth to mouth. But moveable type enabled books to be mass produced at lower costs. In our time, virtually everything is available. Nevertheless, the widespread availability of esoteric teachings has not prevented them from being misunderstood.

For example, New Age teachings mimic esoteric teachings while distorting them out of all recognition. For a while, I followed such teachings, especially the branch that was summarized in the movie The Secret. Called New Thought, the idea was that changing one's thoughts would automatically change the outward circumstances of life. There is no doubt that an improvement in the quality of one's thoughts is beneficial. Unfortunately, however, this won't ``cause" the Universe to mold circumstances to fulfill your desires. Especially, a ten minute affirmation on success cannot override a full day of ``stinkin' thinking'".

They claimed to be teaching what the ``mystics always taught." However, I could not recall mystics trying to use their spirituality to gain wealth, health, love, and the like. That was when I made the decision to actually read those mystics for myself. Although New Thought had the words of the mystics, it did not translate into an understanding.

This is one of the affirmations we were taught: ``I am Perfect, Whole, and Complete." That is worth exploring. The fundamental misunderstanding is the failure to grasp the distinction between the Self (or Person) and the Ego (or Individual). So the affirmations may be correct when applied to the Self, but not to the Individual.

\paragraph{Perfect}
The Self is perfect but the individual ego certainly is not. Calling the Ego ``prefect" is actually harmful. Alchemical work begins with purgation. One must honestly face up to the shadow of the ego and engage in the process of overcoming negative thoughts, images, and emotions. Repeating affirmations will not accomplish that. Unfortunately, in new age circles that work is avoided because it is considered to be purely negative.

\paragraph{Whole}
The Self is an undivided whole but the individual is full of conflicting elements. The next stage in the alchemical process is to resolve those conflicts. To be Whole means to be a Unity. The individual who aspires to such wholeness will develop a singular will that will remain constant from day to day.

\paragraph{Complete}
The Self is Complete because it does not lack anything. The New Ages recognize that the individual experiences lack. It may be deficient financially, health.

They claim that the opposite of lack is supply. Supply can be replenished by contacting the ``higher self". In a sense that is true. However, one does not ``contact" the higher self, one must ``be" the higher self. That is the final stage of alchemical transformation.

\paragraph{Summary}
The new ager who believes that the individual as such is perfect, whole, and complete, will fail to develop spiritually in a fully healthy way.

Nevertheless, it may be useful to affirm to oneself: ``I am perfect, whole, and complete." That will serve as a reminder that you, as a Self, have those qualities. Understood properly, it will motivate you to complete the alchemical transformation in your own life. The Universe will not do it for you.



\flrightit{Posted on 2023-03-23 by Cologero }
